\subsection{Overall Progress}
% All statistics in this section use best-of-three (max per CVE-model pair across 3 runs) -> per-pair average.
Table~\ref{tab:model-results} reports the best-of-three results for each model across 30 CVEs. \texttt{glm-5.2} achieves the highest average overall score (64.6/100), followed by \texttt{mimo-v2.5} (52.9) and \texttt{claude-sonnet-4-6} (50.0). While real-target success remains rare---only 11 out of 150 CVE--model pairs achieve near-full P5 credit ($\geq$15/20), of which 8 also achieve near-full P6 credit---the best single run reaches 99.0/100, demonstrating that full autonomous reproduction is achievable under favorable conditions.


Specifically, five CVEs achieve full or near-full reproduction (overall $\geq$90), each enabled by a distinct rehosting strategy. Same-architecture \textit{chroot} accounts for 5 of 12 successful runs on CVE-2023-26315 (AArch64), where the host and target share the same instruction set and no emulation is needed. QEMU system mode with a bootable OS image enables 3 runs on CVE-2025-6443, whose CHR (Cloud Hosted Router) image boots as a complete system without requiring filesystem extraction. FastCGI socket injection enables 3 runs on CVE-2023-44418, where the agent bypasses the lighttpd frontend and sends protocol frames directly to \texttt{prog.cgi}. The remaining run on CVE-2025-60690 uses chroot with cross-architecture QEMU user-mode emulation on a MIPS target, augmented with LD\_PRELOAD shims for fork and CGI handling. CVE-2024-5293 achieves near-full reproduction via Unicorn emulation of the real vulnerable function bytes. In contrast, devices whose firmware is unavailable or requires complex cross-architecture emulation with proprietary NVRAM dependencies almost universally fail at P5.


Notably, the leading model, \texttt{glm-5.2}, wins not through raw reasoning but through task-appropriate behavior. It averages 11.5/15 on P2 (Firmware Acquisition) under best-of-three scoring, well above \texttt{claude-sonnet-4-6}'s 7.7/15; since P2 gates all downstream phases, this gap cascades---in 5 out of 30 CVE--model pairs, \texttt{claude-sonnet-4-6}'s best plan exceeds 50 yet its best task stalls below 15, all blocked at P2. \texttt{glm-5.2} also achieves the strongest fallback planning (50.8/100 vs.\ 27.5--43.3 for the others), which correlates with deeper pipeline progression. On from-scratch reproduction, persistence in artifact acquisition and fallback-driven recovery matter more than model size.

Finally, we also report the computational cost of each model in our appendix, including per-run token consumption, API cost, and execution time across all 450 runs.
